\documentclass{article}

\usepackage{amsmath}
\usepackage{amssymb}

\newcommand{\ZZ}{\mathbb{Z}}
\newcommand{\RR}{\mathbb{R}}
\newcommand{\CC}{\mathbb{C}}

\begin{document}

{\bf Worksheet \#7; date: 09/19/2018}

{\bf MATH 55 Discrete Mathematics}

\begin{enumerate}
\item {\em (Rosen 4.3.11)} Show that $\log_2 3$ is an irrational number. Recall that an irrational number is a real number $x$ that cannot be written as the ratio of two integers.

\item {\em (Rosen 4.3.17c)} Determine whether the integers in the following set are pairwise relatively prime:
\[
12, 17, 31, 37
\]

\item {\em (Rosen 4.3.25a, f)} What are the greatest common divisors of these pairs of integers?
\begin{enumerate}
\item $3^7 \cdot 5^3 \cdot 7^3$, $2^{11} \cdot e^5 \cdot 5^9$
\setcounter{enumii}{5}
\item $1111$, $0$
\end{enumerate}

\item {\em (Rosen 4.3.33e, f; modified)} Use Euclidean algorithm to find
\begin{enumerate}
\setcounter{enumii}{4}
\item $\gcd(1000, 5040)$
\item $\gcd(9888, 6060)$
\end{enumerate}
Write the gcd as a linear combination of the two numbers as well.

\item {\em (Rosen 4.3.50)} Show that if $a$, $b$ and $m$ are integers such that $m \ge 2$ and $a \equiv b \pmod{m}$, then $\gcd(a, m) = \gcd(b, m)$.

\item {\em (Rosen 4.3.55)} Adapt the proof in the text that there are infinitely many primes to prove that there are infinitely many primes of the form $4k+3$, where $k$ is a nonnegative integer. {\em Hint:} Suppose that there are only finitely many such primes $q_1, q_2, \ldots, q_n$, and consider the number $4q_1 q_2 \cdots q_n - 1$.
\end{enumerate}

\end{document}